% Section 9.3, from lectures/L09/README.md: what you should be able to do, the questions to test
% yourself with, and the next lecture. The closing paragraph is from the README's reference list.

\needspace{10\baselineskip}
\section{Sammanfattning}
\label{c9:sec:summary}

\needspace{6\baselineskip}
\subsection*{Det här ska du kunna nu}
Efter det här kapitlet ska du kunna:
\begin{itemize}
  \item Rita en flödesplan för ett program med beslut och loopar, och översätta den till
    assembler.
  \item Förklara vad \code{cp} och \code{cpi} gör med flaggorna, och välja rätt villkorligt hopp
    för ett villkor.
  \item Skriva om-annars, och förklara varför den första grenen slutar med ett \code{rjmp}.
  \item Välja mellan hopp för tal utan tecken och tal med tecken.
  \item Skriva en loop som går ett bestämt antal varv, och säga hur många varv den går för
    startvärdet 0.
  \item Räkna ut hur många klockcykler en fördröjning tar, och välja konstanter för en önskad tid.
  \item Mäta en tid i simulatorn med brytpunkter och stoppuret, och förklara en skillnad mot
    räkningen.
  \item Skriva ett flerval med en gemensam utgång.
\end{itemize}

\needspace{6\baselineskip}
\subsection*{Frågor att testa dig själv med}
\begin{enumerate}
  \item \code{cp r16, r17} följt av \code{brlo}: hoppar programmet om \code{r16} är mindre än
    \code{r17}, eller tvärtom?
  \item Varför behövs ingen \code{cpi} före \code{brne} i en loop som räknar ned med \code{dec}?
  \item Vad händer om \code{rjmp} efter den första grenen i ett om-annars saknas?
  \item En loop räknar ned från startvärdet 0. Hur många varv går den?
  \item Fördröjningen \code{ldi r18, 200} / \code{dec r18} / \code{brne} tar 600 cykler, inte 599
    och inte 800. Varför?
  \item Varför går det inte att göra en fördröjning på en sekund med två loopar i varandra?
  \item Hur mäter du i simulatorn hur lång tid en del av programmet tar?
\end{enumerate}

Referensbladet i bilaga~\ref{app:avr} har alla hopp och varje instruktions antal klockcykler, och
\secref{studio:4-4} visar hur brytpunkterna och stoppuret används i Microchip Studio.

\needspace{6\baselineskip}
\subsection*{Nästa kapitel}
Kapitel~\ref{c10:ch} handlar om inporten, subrutiner och stacken:
\begin{itemize}
  \item ASCII-koder: att göra om ett tal till tecken.
  \item Tal med 16 bitar, och en fördröjning på en hel sekund.
  \item Inporten: att läsa knappar med \code{PIND}, och varför en nedtryckt knapp läses som noll.
  \item Subrutiner: att skriva fördröjningen en gång och anropa den många gånger.
  \item Stacken: var returadressen hamnar, och hur register sparas med \code{push} och \code{pop}.
\end{itemize}
